 % ── Ⅱ. 본 론 ────────────────────────────────────────────
  \chapter{본\quad 론}
  \section{연구방법 및 세부 내용}
  본론은 1~3개 내외의 장으로 구성하며, 연구내용, 연구방법 및 연구결과를 기술한다. 
  서론에서 정의한 문제(연구질문)에 대한 답을 찾아가는 과정을 논리적이고 일목요연하게 기술한다.

  \subsection{그림 및 표 양식 검증}
  표 내부 텍스트 고딕 처리 및 캡션 왼쪽 정렬 지침이 완벽하게 가동되는 예제 데이터 세트입니다.

  \begin{table}[htbp]
    \caption{Title of the Table}
    \label{tab:sample_table}
    \centering
    \setlength{\tabcolsep}{10pt}
    \setlength{\extrarowheight}{6pt}
    \begin{tabular}{l|l|c|c|c}
      \hline
      \multicolumn{2}{l|}{} & \makecell{\textbf{Diameter}\\\textbf{(mm)}} & \makecell{\textbf{Sectional Area}\\\textbf{(mm\textsuperscript{2})}} & \makecell{\textbf{Effective Stress}\\\textbf{(MPa)}} \\ \hline
      Rebar  & Vertical (\#14) & 43.0 & 1452 & -- \\[4pt] \hline
      Tendon & Vertical        & 86.5 & 5877 & 855.8 \\[4pt] \hline
    \end{tabular}
  \end{table}